\documentclass[12pt]{ctexart}
\usepackage[a4paper,margin=2.5cm]{geometry}
\usepackage{amsmath,amssymb}
\usepackage{graphicx}
\usepackage{booktabs}
\usepackage{longtable}
\usepackage{hyperref}
\title{题目待填写}
\author{}
\date{}
\begin{document}
\maketitle
\begin{abstract}
摘要应概述问题、模型、算法、结果和结论。电子版论文第一页即为摘要页。
\end{abstract}
\noindent\textbf{关键词：} 数学建模；模型；算法；验证
\section{问题重述}
重述赛题问题和需要回答的子问题。
\section{模型假设}
列出必要、合理且可检验的假设。
\section{符号说明}
\begin{longtable}{lll}
\toprule
符号 & 含义 & 单位 \\
\midrule
\endhead
\bottomrule
\end{longtable}
\section{问题分析}
分析问题结构、数据条件和建模路线。
\section{模型建立}
给出模型、变量、目标函数和约束条件。
\section{模型求解}
说明算法、参数、计算过程和复现命令。
\section{结果与解释}
插入代码生成的结果图表，并解释其含义。
\section{模型检验}
给出误差分析、合理性验证或对比验证。
\section{灵敏度分析}
分析关键参数变化对结果的影响。
\section{模型评价}
说明模型优点、局限和改进方向。
\section{结论}
归纳主要结论和答题结果。
\begin{thebibliography}{9}
\bibitem{cumcm-format} 全国大学生数学建模竞赛组委会. 全国大学生数学建模竞赛论文格式规范（2026年修订稿）.
\end{thebibliography}
\section*{附录：支撑材料文件列表}
列出支撑材料文件、可运行源程序、数据文件、较大篇幅中间结果图表和 AI 使用详情或未使用声明。
\end{document}
